\chapter{BRAF Genetic Testing}

\section*{What Is This Test?}
BRAF genetic testing detects somatic (acquired) mutations in the BRAF gene, a proto-oncogene located on chromosome 7q34 that encodes the B-Raf serine/threonine protein kinase, a key component of the RAS-RAF-MEK-ERK (MAP kinase/ERK) signal transduction pathway. The most common BRAF mutation is the V600E point mutation (a substitution of valine by glutamic acid at codon 600 in exon 15), which results in constitutive activation of the BRAF kinase and downstream signaling, driving unregulated cell proliferation independent of growth factor stimulation. BRAF V600E is found in approximately 50--60\% of cutaneous melanomas, 10--15\% of colorectal adenocarcinomas (sporadic mismatch repair-proficient), 45--70\% of papillary thyroid carcinomas, nearly 100\% of hairy cell leukemias, 50--60\% of Langerhans cell histiocytoses, and a subset of gliomas, non-small cell lung cancers (1--2\%), and ovarian cancers. BRAF V600E testing is a companion diagnostic for targeted therapy with BRAF inhibitors (vemurafenib, dabrafenib, encorafenib) as monotherapy or in combination with MEK inhibitors (cobimetinib, trametinib, binimetinib). In colorectal cancer, BRAF V600E mutation is a strong negative prognostic marker and predicts lack of response to anti-EGFR monoclonal antibodies (cetuximab, panitumumab). Testing is performed by PCR-based methods (allele-specific PCR, real-time PCR, droplet digital PCR), Sanger sequencing, or next-generation sequencing (NGS) on formalin-fixed paraffin-embedded tumor tissue.

\section*{Alternative Names}
BRAF V600E; BRAF Mutation Analysis; BRAF V600 Mutation; BRAF Codon 600 Mutation; B-Raf Mutation Testing; BRAF Genotyping

\section*{Principle}
BRAF mutation testing is performed on DNA extracted from formalin-fixed paraffin-embedded (FFPE) tumor tissue sections. The most common assay is real-time PCR using allele-specific probes (TaqMan or Scorpion-ARMS technology) designed to detect the specific BRAF V600E (c.1799T>A) and V600K (c.1798\_1799GT>AA) mutations with high sensitivity (detection limit ~1--5\% variant allele frequency in a background of wild-type DNA). The cobas 4800 BRAF V600 Mutation Test (Roche, FDA-approved companion diagnostic) uses real-time PCR on the cobas z 480 analyzer. Droplet digital PCR (ddPCR) provides absolute quantification of the V600E mutation with detection sensitivity down to 0.01\%. Next-generation sequencing (NGS) of tumor DNA using targeted amplicon or hybrid-capture panels can simultaneously assess BRAF and multiple other actionable mutations (RAS, EGFR, PIK3CA, etc.) with sensitivity of approximately 5\% variant allele frequency. Sanger sequencing has lower sensitivity (~15--20\%) and is no longer recommended as the sole method for BRAF testing. For melanomas, testing should include not only V600E but also V600K, V600D, and V600R mutations (rare but targetable).

\section*{History}
The BRAF gene was first identified as a viral oncogene homolog by Ikawa et al. in 1988, but its significance to human cancer was established by Davies et al. in 2002 in the landmark Cancer Genome Project at the Wellcome Trust Sanger Institute, which reported that BRAF is mutated in approximately 66\% of malignant melanomas and at lower frequencies in other cancers. The BRAF V600E mutation was shown to constitutively activate the MAP kinase pathway. In 2011, vemurafenib (PLX4032, Zelboraf) was approved by the FDA for BRAF V600E-mutated metastatic melanoma after the BRIM-3 trial demonstrated a significant improvement in progression-free and overall survival compared to dacarbazine. Dabrafenib (Tafinlar) and encorafenib (Braftovi) were subsequently approved. The addition of a MEK inhibitor (trametinib, cobimetinib, binimetinib) was shown to further improve progression-free survival and reduce toxicity (cutaneous squamous cell carcinomas from paradoxical MAPK activation). In colorectal cancer, the understanding of BRAF V600E as a strong negative prognostic marker and predictor of anti-EGFR resistance emerged from multiple studies between 2008--2013.

\section*{Why This Test Is Ordered}
Determination of BRAF V600 mutation status in metastatic melanoma to guide targeted therapy with a BRAF inhibitor (vemurafenib, dabrafenib, encorafenib) plus a MEK inhibitor (cobimetinib, trametinib, binimetinib). Patients with BRAF V600E/K-mutated melanoma are candidates for first-line BRAF/MEK inhibitor combination therapy. In metastatic colorectal cancer, BRAF V600E testing is performed simultaneously with RAS (KRAS/NRAS) mutation testing and MSI/dMMR testing to guide treatment decisions: BRAF V600E predicts lack of response to anti-EGFR antibodies (cetuximab, panitumumab) in RAS wild-type tumors. BRAF V600E in colorectal cancer is a negative prognostic marker. In papillary thyroid carcinoma, BRAF V600E status is used for risk stratification and may guide the extent of surgery. In hairy cell leukemia, BRAF V600E confirms the diagnosis and guides treatment with BRAF inhibitors (vemurafenib). In Langerhans cell histiocytosis and Erdheim-Chester disease, BRAF V600E guides vemurafenib therapy. In non-small cell lung cancer, BRAF V600E (found in 1--2\%) guides treatment with dabrafenib plus trametinib. In low-grade gliomas, particularly pilocytic astrocytomas and gangliogliomas, BRAF V600E status guides targeted therapy.

\section*{Is This a Primary or Secondary Test}
This is a primary companion diagnostic test required for the selection of BRAF-targeted therapy in melanoma, and a primary predictive/prognostic test in colorectal cancer.

\section*{How This Test Is Performed}
BRAF testing is performed on tumor tissue obtained by biopsy or surgical excision. The tissue is fixed in formalin, embedded in paraffin, sectioned, and the tumor area is identified on an H\&E-stained slide. Tumor-enriched areas are macrodissected or microdissected to ensure $\geq$20\% tumor cell content (to achieve adequate sensitivity). DNA is extracted from the tumor tissue and analyzed by the PCR or NGS method. Results are typically available in 3--10 days.

\section*{What Preparation Is Needed}
No patient preparation is required. The key pre-analytical requirement is adequate tissue fixation: 10\% neutral buffered formalin for 6--48 hours. Acid-based decalcification (for bone biopsies) degrades DNA and can cause false-negative results. The minimum tumor cell content for reliable detection is approximately 20\% for standard PCR methods and 10--20\% for NGS.

\section*{Contraindications and Precautions}
No absolute contraindications. In colorectal cancer, BRAF V600E mutation is mutually exclusive with KRAS and NRAS mutations in most patients due to functional redundancy in the MAP kinase signaling pathway. However, testing for BRAF and RAS should be performed concurrently because RAS testing is always required in metastatic colorectal cancer. A BRAF V600E mutation found in a colorectal cancer with deficient mismatch repair (dMMR/MSI-H) should raise the possibility of Lynch syndrome, particularly if the patient is young. However, the BRAF V600E mutation itself strongly suggests that the tumor is sporadic (methylation of the MLH1 promoter) rather than Lynch syndrome, because BRAF V600E is nearly exclusively found in sporadic MSI-H colorectal cancer rather than Lynch-associated tumors. If BRAF V600E is present in a dMMR tumor, germline genetic testing for Lynch syndrome may not be required. The cobas 4800 BRAF V600 test is FDA-approved only for melanoma; its use in colorectal cancer is off-label but clinically standard. NGS-based comprehensive genomic profiling can simultaneously assess BRAF, RAS, MSI, TMB, NTRK, HER2, and other actionable alterations. BRAF V600E testing on circulating tumor DNA (ctDNA, liquid biopsy) has lower sensitivity (70--80\%) compared to tissue testing and should only be used when tissue is insufficient.

\section*{Risks}
Risks are related to the biopsy procedure. The greater risk is a false-negative BRAF result from insufficient tumor content, poor DNA quality, or testing that fails to detect non-V600E BRAF mutations that may still be targetable.

\section*{What Happens After the Test}
A positive BRAF V600E result in metastatic melanoma confirms eligibility for BRAF/MEK inhibitor combination therapy, which is the preferred first-line systemic therapy for BRAF-mutated melanoma per NCCN guidelines. A positive BRAF V600E result in colorectal cancer predicts a poor response to anti-EGFR antibodies and, combined with MSI-H/dMMR, defines a patient subset with poor prognosis but high likelihood of response to immunotherapy (pembrolizumab or nivolumab plus ipilimumab, independent of BRAF status). A negative BRAF result should be confirmed to include V600K and other V600 variants if the clinical scenario is melanoma (approximate testing for non-V600E mutations).

\section*{Understanding Results}

\subsection*{BRAF V600E Mutation Negative}
Wild-type BRAF at codon 600. In melanoma, the patient is not eligible for BRAF/MEK inhibitor therapy; first-line treatment is typically immunotherapy (anti-PD-1 monotherapy or nivolumab plus ipilimumab). In colorectal cancer, the absence of BRAF V600E with RAS wild-type status indicates that the patient is eligible for anti-EGFR antibody therapy (cetuximab or panitumumab) as monotherapy or in combination with chemotherapy. However, rare non-V600 BRAF mutations (Class 2 and Class 3 mutations) may be present; these are not targetable by approved BRAF inhibitors but confer variable prognosis. In papillary thyroid carcinoma, absence of BRAF V600E is a favorable prognostic factor that correlates with a lower probability of extrathyroidal extension, lymph node metastases, and recurrence.

\subsection*{BRAF V600E Mutation Positive}
In metastatic melanoma, a positive BRAF V600E result confirms eligibility for BRAF/MEK inhibitor combination therapy, which produces response rates of 60--70\% and median progression-free survival of 11--14 months. In metastatic colorectal cancer, BRAF V600E is present in 10--15\% of patients and is associated with: (1) a very poor prognosis (median overall survival of 12--18 months versus 30+ months for BRAF/RAS wild-type), (2) resistance to anti-EGFR antibodies (cetuximab, panitumumab), regardless of RAS status, and (3) an increased likelihood of MSI-H (approximately 20\% of BRAF-mutated CRCs are MSI-H, and these patients respond well to immunotherapy). BRAF V600E in colorectal cancer with proficient MMR: the BEACON CRC trial demonstrated that encorafenib plus cetuximab (with or without binimetinib) improves overall survival compared to standard chemotherapy. In papillary thyroid carcinoma, BRAF V600E is an adverse prognostic factor. In hairy cell leukemia, BRAF V600E is present in virtually 100\% of cases and has become a diagnostic criterion.

\section*{Normal Results}
BRAF V600E mutation NOT DETECTED (wild-type). The normal BRAF gene encodes the B-Raf protein that functions normally in the MAP kinase pathway, activated by RAS-GTP and subsequently activating MEK, which activates ERK, which translocates to the nucleus to regulate gene expression. BRAF mutation testing is a qualitative test; no quantitative ``normal range'' exists. The result is reported as ``mutation detected'' or ``mutation not detected,'' with specification of the specific mutation identified (e.g., BRAF V600E c.1799T>A).

\section*{Follow-up Tests}
RAS (KRAS and NRAS) mutational testing (exons 2, 3, 4) in colorectal cancer, mandatory because BRAF and RAS status jointly determine anti-EGFR eligibility. Microsatellite instability (MSI) testing or mismatch repair (MMR) IHC (MLH1, MSH2, MSH6, PMS2), essential in all colorectal cancers for Lynch syndrome screening and immunotherapy selection. Comprehensive genomic profiling (NGS panel) to identify additional actionable mutations (NTRK fusions, HER2 amplification, PIK3CA mutations, RET fusions, ALK/ROS1 fusions). Circulating tumor DNA (ctDNA) testing (liquid biopsy) for BRAF V600E when tissue is insufficient or to monitor treatment response and acquired resistance. Serial BRAF testing from ctDNA to detect emergence of resistance mutations during BRAF inhibitor therapy. MLH1 promoter methylation testing: if a colorectal cancer is dMMR and BRAF V600E is positive, MLH1 promoter hypermethylation testing confirms sporadic etiology, potentially avoiding unnecessary germline testing for Lynch syndrome. Thyroid ultrasound with cervical lymph node mapping if BRAF V600E-positive papillary thyroid carcinoma is diagnosed.

\section*{References}
Davies H, Bignell GR, Cox C, et al. Mutations of the BRAF gene in human cancer. Nature. 2002;417(6892):949--954. Chapman PB, Hauschild A, Robert C, et al. Improved survival with vemurafenib in melanoma with BRAF V600E mutation (BRIM-3). N Engl J Med. 2011;364(26):2507--2516. Kopetz S, Grothey A, Yaeger R, et al. Encorafenib, binimetinib, and cetuximab in BRAF V600E-mutated colorectal cancer (BEACON CRC). N Engl J Med. 2019;381(17):1632--1643. Long GV, Stroyakovskiy D, Gogas H, et al. Dabrafenib and trametinib versus dabrafenib and placebo for Val600 BRAF-mutant melanoma. Lancet. 2015;386(9992):444--451. NCCN Guidelines: Melanoma: Cutaneous (V2.2024); Colon Cancer (V2.2024); Thyroid Carcinoma (V3.2024).

\clearpage
